% !TEX program = pdflatex
% =====================================================================
%  Informe de las semanas 13 y 14
%  Persistencia, servicios y API del corte vertical (RF-05 y RF-34)
%  Arquitectura de Sistemas I
%  Compilar en Overleaf con pdfLaTeX
% =====================================================================
\documentclass[11pt,a4paper]{article}

\usepackage[utf8]{inputenc}
\usepackage[T1]{fontenc}
\usepackage{newtxtext,newtxmath}
\usepackage{newtxtt}
\usepackage[spanish,es-noshorthands]{babel}
\usepackage[a4paper,top=2.4cm,bottom=2.2cm,left=2.5cm,right=2.5cm,headheight=14pt,headsep=6mm,footskip=10mm]{geometry}
\usepackage{tabularx}
\usepackage{array}
\usepackage{longtable}
\usepackage{booktabs}
\usepackage{titlesec}
\usepackage{fancyhdr}
\usepackage{lastpage}
\usepackage{enumitem}
\usepackage{microtype}
\usepackage{listings}
\usepackage{hyperref}

% ---------------------------------------------------------------------
% Encabezado y pie
% ---------------------------------------------------------------------
\pagestyle{fancy}
\fancyhf{}
\renewcommand{\headrulewidth}{0.4pt}
\renewcommand{\footrulewidth}{0pt}
\fancyhead[L]{\small\scshape Arquitectura de Sistemas I}
\fancyhead[R]{\small\scshape Semanas 13 y 14}
\fancyfoot[C]{\small Página \thepage\ de \pageref{LastPage}}

% ---------------------------------------------------------------------
% Títulos
% ---------------------------------------------------------------------
\titleformat{\section}{\large\bfseries}{\thesection.}{0.6em}{}
\titlespacing*{\section}{0pt}{3.0ex plus .5ex minus .2ex}{1.2ex}
\titleformat{\subsection}{\normalsize\bfseries}{\thesubsection.}{0.6em}{}
\titlespacing*{\subsection}{0pt}{2.2ex plus .4ex minus .2ex}{0.9ex}

% ---------------------------------------------------------------------
% Párrafos, listas y tablas
% ---------------------------------------------------------------------
\setlength{\parindent}{0pt}
\setlength{\parskip}{0.45em}
\setlength{\tabcolsep}{4pt}
\renewcommand{\arraystretch}{1.1}
\setlength{\emergencystretch}{2em}
\raggedbottom

\setlist[itemize]{leftmargin=1.4em,itemsep=1pt,topsep=3pt,parsep=0pt}

\newcolumntype{L}[1]{>{\raggedright\arraybackslash}p{#1}}
\newcolumntype{A}[1]{>{\raggedright\arraybackslash}p{#1}}
\newcolumntype{Y}{>{\raggedright\arraybackslash}X}
\newcolumntype{V}[1]{>{\hsize=#1\hsize\raggedright\arraybackslash}X}
\newcolumntype{Z}[1]{>{\hsize=#1\hsize\centering\arraybackslash}X}

\newcommand{\code}[1]{\texttt{#1}}
\newcommand{\thead}[1]{\textbf{#1}}
\newcommand{\pass}[1]{\textbf{#1}}
\newcommand{\pend}[1]{\textit{#1}}

\lstset{
  basicstyle=\ttfamily\scriptsize,
  frame=lines,
  framerule=0.4pt,
  columns=fullflexible,
  keepspaces=true,
  showstringspaces=false,
  breaklines=true,
  extendedchars=true,
  inputencoding=utf8,
  literate={á}{{\'a}}1 {é}{{\'e}}1 {í}{{\'i}}1 {ó}{{\'o}}1 {ú}{{\'u}}1
           {Á}{{\'A}}1 {É}{{\'E}}1 {Í}{{\'I}}1 {Ó}{{\'O}}1 {Ú}{{\'U}}1
           {ñ}{{\~n}}1 {Ñ}{{\~N}}1 {ü}{{\"u}}1 {Ü}{{\"U}}1
           {¿}{{?`}}1 {¡}{{!`}}1,
  aboveskip=6pt,
  belowskip=6pt
}

\hypersetup{
  hidelinks,
  pdftitle={Informe de las semanas 13 y 14 - Corte vertical},
  pdfsubject={Persistencia, servicios y API de RF-05 y RF-34},
  pdfauthor={Arquitectura de Sistemas I - Universidad Central}
}

\begin{document}

% =====================================================================
% Portada
% =====================================================================
\begin{titlepage}
\centering
\vspace*{1.0cm}

{\scshape\Large Universidad Central\par}
\vspace{0.25cm}
{\scshape\large Arquitectura de Sistemas I\par}

\vspace{2.2cm}
{\Huge\bfseries Informe de las semanas 13 y 14\par}
\vspace{0.45cm}
{\LARGE Persistencia, servicios y API del corte vertical\par}
\vspace{0.35cm}
{\Large\itshape Proyecto Arquitectura de Sistemas\par}
\vspace{0.15cm}
{\large Adopción responsable de mascotas\par}

\vspace{2.2cm}
{\normalsize
\begin{tabular}{@{}l@{\hspace{6mm}}>{\raggedright\arraybackslash}p{9.2cm}@{}}
\textbf{Proyecto:} & Sistema de Apoyo a la Decisión para la Adopción Responsable de Mascotas mediante Inteligencia Artificial \\[3mm]
\textbf{Docente:} & Elias Buitrago Bolivar \\[3mm]
\textbf{Integrantes:} & Juan Diego Rojas y Adrián Camaño \\[3mm]
\textbf{Fecha:} & 10 de noviembre de 2026 \\[3mm]
\textbf{Estado:} & 12/12 pruebas de persistencia y 39/39 pruebas de API \\
\end{tabular}}

\vfill
{\small Semanas 13 y 14 \textperiodcentered\ 2026\par}
\end{titlepage}

% =====================================================================
% 1. Resumen
% =====================================================================
\section{Resumen de la entrega}

Este documento reúne el trabajo de las semanas 13 y 14 sobre el mismo corte
vertical. La semana 13 construyó la base de datos y la capa de persistencia; la
semana 14 agregó los servicios de aplicación, la API REST, los contratos y la
verificación de los requisitos no funcionales. Las dos semanas conservan los
mismos dos requisitos funcionales priorizados.

{\footnotesize
\begin{tabularx}{\linewidth}{@{}V{0.35}V{1.0}V{1.0}V{1.65}@{}}
\toprule
\thead{Semana} & \thead{Trabajo realizado} & \thead{Resultado} & \thead{Evidencia} \\ \midrule
\textbf{13} & Esquema, reglas de integridad, repositorio y datos de prueba &
12/12 pruebas de persistencia & \code{evidencias/salida\_pruebas.txt} \\
\textbf{14} & Servicios, API FastAPI, contratos, pruebas funcionales y RNF &
39/39 pruebas; 100\,\% de cobertura en servicios &
\code{salida\_pruebas\_semana14.txt}, \code{cobertura.txt},
\code{rnf\_desempeno.txt} y \code{demo\_api.txt} \\
\bottomrule
\end{tabularx}
}

\vspace{1mm}
{\footnotesize
\begin{tabularx}{\linewidth}{@{}V{0.3}V{1.2}V{1.5}@{}}
\toprule
\thead{Orden} & \thead{Requisito funcional} & \thead{Resultado observable} \\ \midrule
1 & \textbf{RF-05 --- Gestión de la ficha de mascota} &
El refugio crea, edita y da de baja una ficha; recibe su identificador y su
estado (\emph{incompleta} o \emph{disponible}). \\
2 & \textbf{RF-34 --- Consulta del catálogo} &
El adoptante lista las mascotas disponibles con filtros por especie, tamaño y
edad, con paginación y orden estable. \\
\bottomrule
\end{tabularx}
}

\textbf{Stack:} Python 3.12.10, FastAPI 0.162.2, Pydantic v2, Uvicorn 0.54.0 y
SQLite 3.49.1. SQLite se usa para reproducir el corte sin instalar un servidor;
la base objetivo del proyecto sigue siendo PostgreSQL (ADR-03/ADR-10).

% =====================================================================
% 2. Contexto
% =====================================================================
\section{Contexto y alcance}

\subsection{El caso de estudio}

El sistema nace en la asignatura \emph{Práctica de Ingeniería}. Los refugios
registran las fichas de los animales que tienen en adopción y los adoptantes
consultan ese catálogo para encontrar una mascota compatible con su hogar. En
\textbf{Arquitectura de Sistemas I} se retoma como caso real para aplicar
decisiones de arquitectura; no se construye el sistema completo, sino un
\textbf{corte vertical}.

Un corte vertical es una porción pequeña del sistema que se construye de punta a
punta. En este proyecto avanza en tres momentos: la semana 13 cubre la base de
datos y la persistencia; la semana 14, los servicios y la API; la semana 3 del
prototipo, la interfaz, la autenticación y el cierre.

\subsection{Los dos requisitos funcionales}

\begin{itemize}
\item \textbf{RF-05.} El refugio crea, edita y da de baja la ficha de una
  mascota. La ficha incluye especie, raza aproximada, edad, tamaño, nivel de
  energía, historial de salud y comportamiento, experiencia requerida, aptitud
  para convivir con menores y otras mascotas, y de 1 a 5 fotografías.
\item \textbf{RF-34.} El adoptante consulta el catálogo de mascotas disponibles
  y lo filtra por especie, tamaño y rango de edad, con paginación y orden por
  fecha de publicación.
\end{itemize}

Se eligieron porque \textbf{RF-05 produce la información y RF-34 la consume}.
Juntos forman el recorrido mínimo del negocio ---publicar una mascota y
consultarla--- y comparten las tablas \emph{refugio}, \emph{mascota} y
\emph{fotografia}.

\subsection{Reglas de negocio del corte}

{\footnotesize
\begin{tabularx}{\linewidth}{@{}V{0.4}V{1.7}V{0.9}@{}}
\toprule
\thead{Regla} & \thead{Enunciado} & \thead{Dónde se aplica} \\ \midrule
\textbf{RN-22} & Una ficha está completa cuando tiene los campos obligatorios y
entre 1 y 5 fotos JPG/PNG; solo entonces pasa a \emph{disponible}. &
Repositorio, triggers y API \\
\textbf{RN-03} & Solo un refugio verificado puede publicar fichas. &
Repositorio y API (ERR-009) \\
\textbf{RN-15} & Una ficha adoptada o dada de baja no aparece en el catálogo. &
Filtro del catálogo y baja lógica \\
\bottomrule
\end{tabularx}
}

\subsection{Fuera del alcance}

Autenticación y roles, interfaz web, almacenamiento binario de fotos, motor de
recomendaciones, solicitudes y decisiones de adopción, seguimiento posadopción,
notificaciones, indicadores, chatbot y reconocimiento de imágenes. La migración
definitiva a PostgreSQL se hará al cerrar el prototipo.

\clearpage

% =====================================================================
% 3. Semana 13
% =====================================================================
\section{Semana 13 --- Base de datos y persistencia}

\subsection{Modelo mínimo de datos}

{\footnotesize
\begin{tabularx}{\linewidth}{@{}V{0.55}V{0.55}V{1.95}V{1.0}V{0.95}@{}}
\toprule
\thead{Entidad} & \thead{Clave} & \thead{Campos principales} & \thead{Relación} & \thead{RF} \\ \midrule
\textbf{refugio} & id (PK) & nombre, localidad, estado\_verificacion
(pendiente, verificado, rechazado) & 1 : N con mascota & RF-05, RF-34 \\
\textbf{mascota} & id (PK) & refugio\_id (FK), nombre, raza, especie,
edad\_meses, tamano, nivel\_energia, experiencia\_requerida,
apta\_con\_menores, apta\_con\_mascotas, historial\_salud,
historial\_comportamiento, estado, fecha\_publicacion &
N : 1 con refugio; 1 : N con fotografia & RF-05, RF-34 \\
\textbf{fotografia} & id (PK) & mascota\_id (FK), archivo, orden (1 a 5) &
N : 1 con mascota; UNIQUE (mascota\_id, orden) & RF-05, RF-34 \\
\bottomrule
\end{tabularx}
}

El esquema se recrea desde cero con \code{persistencia/schema.sql} y
\code{persistencia/seed.sql}. La base de la semana 13 vive en
\code{evidencias/prueba.db} y no se modifica desde la API.

\subsection{Integridad y separación de capas}

\noindent
\begin{minipage}[t]{0.49\linewidth}
\begin{itemize}
\item \textbf{Llaves foráneas} de refugio a mascota y de mascota a fotografia,
  con borrado en cascada de las fotos.
\item \textbf{CHECK} de dominios cerrados: especie, tamaño, energía,
  experiencia y estado; edad entre 0 y 300 meses.
\item \textbf{UNIQUE} en (mascota\_id, orden) y trigger que impide más de cinco
  fotos.
\item \textbf{Triggers de publicación:} una ficha solo pasa a disponible si
  tiene al menos una foto y el refugio está verificado (RN-03 y RN-22).
\end{itemize}
\end{minipage}\hfill
\begin{minipage}[t]{0.49\linewidth}
\begin{itemize}
\item \textbf{Transacciones} con \emph{BEGIN IMMEDIATE} y \emph{ROLLBACK} ante
  error.
\item \textbf{Repositorio único:} todo el SQL está en
  \code{persistencia/repositorio.py}; los servicios y la API no construyen
  consultas.
\item \textbf{Catálogo:} solo fichas disponibles de refugios verificados, sin
  datos personales del adoptante.
\item \textbf{Datos de fotos:} en este corte se guarda la ruta del archivo; el
  binario se incorporará más adelante.
\end{itemize}
\end{minipage}

\subsection{Contrato del repositorio}

{\footnotesize
\begin{tabularx}{\linewidth}{@{}V{1.0}V{0.95}V{1.15}V{0.9}@{}}
\toprule
\thead{Operación} & \thead{Entrada} & \thead{Resultado} & \thead{Uso} \\ \midrule
\code{crear\_ficha} & refugio\_id, ficha de RN-22, 0--5 fotos &
id real; estado incompleta o disponible & RF-05: persistir y devolver el
identificador. \\
\code{obtener\_por\_id} & mascota\_id & ficha con refugio y fotos, o None &
RF-05: verificar el registro; representar la ausencia. \\
\code{actualizar\_ficha} & mascota\_id, refugio\_id, cambios, fotos opcionales &
ficha actualizada y estado recalculado & RF-05: editar y volver a publicar. \\
\code{dar\_de\_baja} & mascota\_id, refugio\_id & baja lógica (estado = baja) &
RF-05: retirar la ficha del catálogo. \\
\code{listar\_catalogo} & especie, tamano, edad\_min, edad\_max, pagina,
por\_pagina & total, pagina, porPagina, mascotas &
RF-34: consulta paginada y filtrada. \\
\bottomrule
\end{tabularx}
}

\subsection{Datos iniciales}

{\footnotesize
\begin{tabularx}{\linewidth}{@{}V{0.8}V{1.2}V{0.9}V{0.75}Z{0.8}V{1.55}@{}}
\toprule
\thead{Mascota} & \thead{Refugio} & \thead{Especie / tamaño} &
\thead{Estado} & \thead{¿Catálogo?} & \thead{Motivo} \\ \midrule
1 --- Luna & 1 --- Patitas (verificado) & perro / grande & disponible &
\textbf{Sí} & Ficha completa. \\
2 --- Milo & 1 --- Patitas (verificado) & gato / pequeño & disponible &
\textbf{Sí} & Ficha completa. \\
3 --- Rocky & 1 --- Patitas (verificado) & perro / mediano & disponible &
\textbf{Sí} & Ficha completa. \\
4 --- Nube & 2 --- Colitas (pendiente) & gato / pequeño & disponible &
\textit{No} & Refugio no verificado. \\
5 --- Toby & 1 --- Patitas (verificado) & perro / pequeño & incompleta &
\textit{No} & Sin fotografía. \\
6 --- Canela & 1 --- Patitas (verificado) & perro / mediano & adoptada &
\textit{No} & Ya adoptada. \\
\bottomrule
\end{tabularx}
}

Catálogo esperado con los datos iniciales: \textbf{Luna (1), Rocky (3), Milo
(2)}, ordenados por \code{fecha\_publicacion} ascendente y, a igual fecha, por
\code{id}.

\subsection{Resultados de persistencia}

Las 12 pruebas de la semana 13 se ejecutan con:

\begin{lstlisting}
cd persistencia
$env:PYTHONUTF8 = "1"
$env:PYTHONIOENCODING = "utf-8"
python prueba_persistencia.py
\end{lstlisting}

{\footnotesize
\begin{longtable}{@{}A{0.05\linewidth}L{0.60\linewidth}L{0.29\linewidth}@{}}
\toprule
\thead{\#} & \thead{Prueba} & \thead{Resultado} \\ \midrule \endfirsthead
\thead{\#} & \thead{Prueba} & \thead{Resultado} \\ \midrule \endhead
\bottomrule
\endlastfoot
1 & Crear ficha en refugio verificado y leerla desde una conexión nueva &
Pasa --- id 7, disponible, 2 fotos. \\
2 & Rechazar ficha de refugio no verificado &
Pasa --- ERR-009; 0 filas nuevas. \\
3 & Rechazar datos inválidos & Pasa --- ERR-004; 0 filas nuevas. \\
4 & Consultar un identificador inexistente & Pasa --- devuelve None. \\
5 & Violar la FK de refugio &
Pasa --- FOREIGN KEY constraint failed. \\
6 & Ficha sin fotos queda incompleta y fuera del catálogo &
Pasa --- estado incompleta. \\
7 & Catálogo con orden, filtros y paginación &
Pasa --- total 4, orden [1, 3, 2, 7]. \\
8 & Actualizar, publicar y dar de baja &
Pasa --- incompleta $\rightarrow$ disponible $\rightarrow$ baja. \\
9 & Rechazar más de cinco fotografías &
Pasa --- repositorio y trigger. \\
10 & Triggers de publicación &
Pasa --- refugio pendiente y ficha sin foto. \\
11 & Catálogo y ficha sin datos personales & Pasa --- 0 campos personales. \\
12 & 20 consultas consecutivas al catálogo &
Pasa --- 0,53 ms promedio; 0,66 ms máximo. \\
\end{longtable}
}

\textbf{Resumen:} 12/12 pruebas superadas. La medición de tiempos de la semana
13 es una línea base con 8 fichas y no sustituye la prueba de carga de RNF-05.
La evidencia completa está en \code{evidencias/salida\_pruebas.txt}.

\clearpage

% =====================================================================
% 4. Semana 14
% =====================================================================
\section{Semana 14 --- Servicios y API}

\subsection{Arquitectura por capas}

\begin{center}
\code{HTTP} $\rightarrow$ \code{API} $\rightarrow$
\code{Servicio de aplicación} $\rightarrow$ \code{Repositorio} $\rightarrow$
\code{SQLite}
\end{center}

{\footnotesize
\begin{tabularx}{\linewidth}{@{}V{0.55}V{1.05}V{1.4}@{}}
\toprule
\thead{Capa} & \thead{Archivo} & \thead{Responsabilidad} \\ \midrule
API & \code{api/main.py} & Recibir HTTP, validar la forma, delegar y traducir
errores a códigos HTTP. No contiene reglas ni SQL. \\
Contratos & \code{api/esquemas.py} & Modelos Pydantic v2 de entrada y salida. \\
Inyección & \code{api/dependencias.py} & Crear el repositorio y los servicios;
la base se toma de \code{DB\_PATH}. \\
Servicios & \code{servicios/} & Coordinar el caso de uso, normalizar la
respuesta y decidir la ausencia de una ficha. \\
Repositorio & \code{persistencia/repositorio.py} & Todo el SQL, las validaciones
de dominio y las transacciones. \\
\bottomrule
\end{tabularx}
}

Esta separación permite probar los servicios con un doble del repositorio, sin
abrir una base real, y mantiene el SQL fuera de la API.

\subsection{Servicios de aplicación}

{\footnotesize
\begin{tabularx}{\linewidth}{@{}V{0.7}V{0.8}V{1.5}@{}}
\toprule
\thead{Servicio} & \thead{Caso de uso} & \thead{Métodos} \\ \midrule
\code{ServicioMascotas} & RF-05 &
\code{crear}, \code{consultar}, \code{actualizar}, \code{dar\_de\_baja} \\
\code{ServicioCatalogo} & RF-34 &
\code{listar} (filtros, paginación y normalización a snake\_case) \\
\bottomrule
\end{tabularx}
}

\code{ServicioMascotas.consultar} traduce la ausencia de una ficha a
\code{FichaNoEncontrada} (ERR-011). \code{ServicioCatalogo} delega los filtros
en el repositorio y convierte \code{porPagina} a \code{por\_pagina} en la
respuesta JSON.

\subsection{Endpoints}

Base local: \url{http://127.0.0.1:8000/api/v1}. Documentación interactiva:
\url{http://127.0.0.1:8000/docs}.

{\footnotesize
\begin{tabularx}{\linewidth}{@{}V{0.35}V{1.35}V{1.2}V{0.9}V{1.2}@{}}
\toprule
\thead{RF} & \thead{Método y ruta} & \thead{Servicio} & \thead{Éxito} &
\thead{Errores posibles} \\ \midrule
RF-05 & \code{POST /mascotas} & \code{ServicioMascotas.crear} &
201 ficha creada & 400 ERR-004, 404 ERR-008, 409 ERR-009 \\
RF-05 & \code{PUT /mascotas/\{id\}} & \code{ServicioMascotas.actualizar} &
200 ficha actualizada & 400 ERR-004, 403 ERR-003, 404 ERR-011 \\
RF-05 & \code{DELETE /mascotas/\{id\}?refugio\_id=\{id\}} &
\code{ServicioMascotas.dar\_de\_baja} & 204 sin cuerpo &
400 ERR-004, 403 ERR-003, 404 ERR-011 \\
RF-34 & \code{GET /mascotas} & \code{ServicioCatalogo.listar} &
200 página del catálogo & 400 ERR-004 \\
\bottomrule
\end{tabularx}
}

Formato único de error:

\begin{lstlisting}
{ "detail": "mensaje legible", "codigo": "ERR-004" }
\end{lstlisting}

{\footnotesize
\begin{tabularx}{\linewidth}{@{}V{0.5}V{0.5}V{2.0}@{}}
\toprule
\thead{Código} & \thead{HTTP} & \thead{Significado} \\ \midrule
ERR-003 & 403 & La ficha no pertenece al refugio indicado. \\
ERR-004 & 400 & Datos de entrada inválidos (tipo, dominio, rango, formato). \\
ERR-006 & 409 & Violación de integridad de la base. \\
ERR-008 & 404 & El refugio no existe. \\
ERR-009 & 409 & El refugio no está verificado. \\
ERR-011 & 404 & La ficha no existe. \\
\bottomrule
\end{tabularx}
}

Las validaciones de forma también responden 400 con ERR-004 (no 422), para
mantener un único contrato de error. Un fallo inesperado responde 500 con
ERR-500 y se registra sin exponer detalles técnicos.

\subsection{Ejemplos de los contratos}

\textbf{Crear una ficha (RF-05).} Cuerpo de \code{POST /mascotas}:

\begin{lstlisting}
{
  "refugio_id": 1,
  "nombre": "Nala",
  "raza": "Mestiza",
  "especie": "perro",
  "edad_meses": 24,
  "tamano": "mediano",
  "nivel_energia": "media",
  "experiencia_requerida": "basica",
  "apta_con_menores": true,
  "apta_con_mascotas": true,
  "historial_salud": "Vacunas al día",
  "historial_comportamiento": "Sociable",
  "fecha_publicacion": "2026-11-10",
  "fotografias": ["nala_01.jpg", "nala_02.jpg"]
}
\end{lstlisting}

Respuesta 201 Created (resumen; la respuesta real incluye todos los campos del
contrato y el id depende de la base; en la base de prueba es 7):

\begin{lstlisting}
{
  "id": 7,
  "refugio_id": 1,
  "nombre": "Nala",
  "especie": "perro",
  "edad_meses": 24,
  "tamano": "mediano",
  "estado": "disponible",
  "fecha_publicacion": "2026-11-10",
  "refugio_nombre": "Refugio Patitas",
  "refugio_localidad": "ENGATIVA",
  "fotografias": [
    { "id": 7, "archivo": "nala_01.jpg", "orden": 1 },
    { "id": 8, "archivo": "nala_02.jpg", "orden": 2 }
  ]
}
\end{lstlisting}

\textbf{Consultar el catálogo (RF-34).} \code{GET /mascotas?especie=perro\&tamano=mediano}
después de crear a Nala en la base de demostración:

\begin{lstlisting}
{
  "total": 2,
  "pagina": 1,
  "por_pagina": 20,
  "mascotas": [
    {
      "id": 3, "nombre": "Rocky", "especie": "perro",
      "edad_meses": 60, "tamano": "mediano",
      "nivel_energia": "alta", "fecha_publicacion": "2026-10-03",
      "refugio_id": 1, "refugio_nombre": "Refugio Patitas",
      "refugio_localidad": "ENGATIVA", "foto_principal": "rocky_01.jpg"
    },
    {
      "id": 7, "nombre": "Nala", "especie": "perro",
      "edad_meses": 24, "tamano": "mediano",
      "nivel_energia": "media", "fecha_publicacion": "2026-11-10",
      "refugio_id": 1, "refugio_nombre": "Refugio Patitas",
      "refugio_localidad": "ENGATIVA", "foto_principal": "nala_01.jpg"
    }
  ]
}
\end{lstlisting}

\textbf{Rechazo por refugio no verificado (RF-05):}

\begin{lstlisting}
{
  "detail": "el refugio 2 no está verificado y no puede publicar fichas",
  "codigo": "ERR-009"
}
\end{lstlisting}

\textbf{Filtro válido sin coincidencias (RF-34):} responde 200 con
\code{total: 0} y \code{mascotas: []}. \textbf{Filtro inválido:} responde 400
con ERR-004. El detalle completo de cada contrato, con todos los campos y
condiciones límite, está en \code{docs/contratos\_api.md}.

\subsection{Decisiones de diseño}

\begin{itemize}
\item \textbf{Actor temporal.} La autenticación y los roles se integran en la
  semana 3. Mientras tanto, el refugio que actúa se envía como
  \code{refugio\_id} en el cuerpo (POST/PUT) o en la consulta (DELETE).
\item \textbf{Validación en dos niveles.} Pydantic valida la forma y el
  repositorio valida el dominio; ambos errores se traducen a ERR-004.
\item \textbf{Base configurable.} La ruta de SQLite se lee de \code{DB\_PATH} en
  cada petición, lo que permite aislar las pruebas sin tocar la base de
  desarrollo.
\item \textbf{CORS preparado.} Los orígenes por defecto son
  \code{http://localhost:5173}, \code{http://127.0.0.1:5173} y
  \code{http://localhost:3000}, configurables con \code{CORS\_ORIGINS}.
\item \textbf{Baja lógica.} DELETE no borra la fila: cambia el estado a baja
  para conservar la trazabilidad y cumplir RN-15.
\item \textbf{Sin datos personales.} El catálogo y la ficha pública solo
  exponen datos de la mascota y del refugio (nombre y localidad), no del
  adoptante.
\end{itemize}

\subsection{Pruebas funcionales}

\begin{lstlisting}
python -m pytest -v
\end{lstlisting}

Resultado: \textbf{39/39 pruebas pasan} en aproximadamente 1,6 s
(\code{evidencias/salida\_pruebas\_semana14.txt}). Aparece un aviso de
deprecación de Starlette sobre \code{httpx}; no afecta los resultados.

{\footnotesize
\begin{longtable}{@{}A{0.06\linewidth}L{0.12\linewidth}L{0.23\linewidth}L{0.29\linewidth}L{0.18\linewidth}@{}}
\toprule
\thead{RF} & \thead{Escenario} & \thead{Acción} & \thead{Resultado obtenido} &
\thead{Evidencia} \\ \midrule \endfirsthead
\thead{RF} & \thead{Escenario} & \thead{Acción} & \thead{Resultado obtenido} &
\thead{Evidencia} \\ \midrule \endhead
\bottomrule
\endlastfoot
RF-05 & Válido & POST con refugio verificado y 2 fotos &
201, id 7, disponible, 2 fotos; verificado en SQLite desde otra conexión &
\code{test\_rf05\_api.py} \\
RF-05 & Rechazo & POST con refugio no verificado &
409 ERR-009; 0 filas nuevas & \code{test\_rf05\_api.py} \\
RF-05 & Entrada & POST con especie inválida &
400 ERR-004; 0 filas nuevas & \code{test\_rf05\_api.py} \\
RF-05 & Edición & PUT de ficha incompleta con una foto &
200; pasa a disponible y aparece en el catálogo & \code{test\_rf05\_api.py} \\
RF-05 & Baja & DELETE de una ficha propia &
204; estado baja y ya no aparece en el catálogo & \code{test\_rf05\_api.py} \\
RF-34 & Válido & GET del catálogo inicial &
200, total 3, orden [1, 3, 2] & \code{test\_rf34\_api.py} \\
RF-34 & Filtros & Especie, tamaño y rango de edad &
[1, 3], [2] y [2] & \code{test\_rf34\_api.py} \\
RF-34 & Paginación & por\_pagina = 2 &
Página 1 [1, 3]; página 2 [2]; total 3 & \code{test\_rf34\_api.py} \\
RF-34 & Ausencia & Filtro válido sin coincidencias &
200, total 0, lista vacía & \code{test\_rf34\_api.py} \\
RF-34 & Entrada & especie = conejo &
400 ERR-004 & \code{test\_rf34\_api.py} \\
\end{longtable}
}

\textbf{Distribución de las 39 pruebas:}

{\footnotesize
\begin{tabularx}{\linewidth}{@{}V{1.4}V{1.1}V{0.5}@{}}
\toprule
\thead{Archivo} & \thead{Propósito} & \thead{Casos} \\ \midrule
\code{pruebas/test\_rf05\_api.py} & Funcionales de RF-05 & 8 \\
\code{pruebas/test\_rf34\_api.py} & Funcionales de RF-34 y RNF-04 & 8 \\
\code{pruebas/test\_rnf10\_integridad.py} & Matriz de entradas inválidas
(RNF-10) & 14 \\
\code{pruebas/test\_servicios\_aislados.py} & Servicios con doble del
repositorio (RNF-08) & 9 \\
\bottomrule
\end{tabularx}
}

\subsection{Demostración HTTP real}

\begin{lstlisting}
python -m pruebas.demo_api
\end{lstlisting}

El script levanta Uvicorn en un hilo, usa un cliente \code{httpx} real y genera
\code{evidencias/demo\_api.txt} con las peticiones y respuestas completas:
crear (RF-05), consultar catálogo (RF-34), actualizar (RF-05), dar de baja
(RF-05), verificar que la ficha dada de baja sale del catálogo, catálogo vacío
y filtro inválido. El resultado final es \textbf{CUMPLE}.

\clearpage

% =====================================================================
% 5. RNF
% =====================================================================
\section{Verificación de requisitos no funcionales}

{\footnotesize
\begin{longtable}{@{}A{0.12\linewidth}L{0.18\linewidth}L{0.21\linewidth}L{0.19\linewidth}L{0.20\linewidth}@{}}
\toprule
\thead{RF / RNF} & \thead{Calidad ISO/IEC 25010:2023} & \thead{Prueba y meta} &
\thead{Resultado} & \thead{Evidencia} \\ \midrule \endfirsthead
\thead{RF / RNF} & \thead{Calidad ISO/IEC 25010:2023} & \thead{Prueba y meta} &
\thead{Resultado} & \thead{Evidencia} \\ \midrule \endhead
\bottomrule
\endlastfoot
RF-05, RF-34 / \textbf{RNF-05} & Eficiencia de desempeño / comportamiento
temporal & 1.000 fichas; 20 consultas consecutivas por HTTP; cada una $\leq$
1 s & \textbf{20/20}; mínimo 2,62 ms; promedio 2,78 ms; p95 3,01 ms; máximo
3,05 ms & \code{evidencias/rnf\_desempeno.txt} \\
RF-05, RF-34 / \textbf{RNF-08} & Mantenibilidad / capacidad para ser probado &
Doble del repositorio, sin base real; cobertura $\geq$ 70\,\% en servicios &
\textbf{9/9} pruebas aisladas; \textbf{100\,\%} de cobertura (33 sentencias) &
\code{test\_servicios\_aislados.py}, \code{cobertura.txt} \\
RF-05, RF-34 / \textbf{RNF-10} & Seguridad / integridad & Matriz de entradas
inválidas; 0 filas persistidas & \textbf{14/14} casos rechazados; 0 filas
nuevas o modificadas & \code{test\_rnf10\_integridad.py},
\code{salida\_pruebas\_semana14.txt} \\
RF-05, RF-34 / \textbf{RNF-04} & Seguridad / confidencialidad & Revisar las
claves del catálogo y de la ficha pública & \textbf{0} campos personales &
\code{test\_rf34\_api.py}, \code{demo\_api.txt} \\
RF-34 / \textbf{RNF-07} & Flexibilidad / escalabilidad & 2.000 fichas y 5.000
adoptantes; p95 $\leq$ 1 s & \pend{Pendiente} para la semana 3 & --- \\
\end{longtable}
}

\subsection{Lectura honesta de RNF-05}

La meta de la plantilla incluye 500 fichas y 50 usuarios concurrentes. En esta
entrega se usó un volumen mayor (1.000 fichas), pero la medición es
\textbf{secuencial}, no concurrente. Por eso el resultado se declara
\textbf{cumplido en lo secuencial} y la prueba de 50 usuarios queda para la
semana 3. No se rebajó el umbral: cada consulta se exigió $\leq$ 1 s.

\subsection{Comandos de verificación}

\begin{lstlisting}
# RNF-08: cobertura de los servicios
python -m coverage run -m pytest -q
python -m coverage report -m

# RNF-05: 1.000 fichas y 20 consultas por HTTP real
python -m pruebas.rnf_desempeno
\end{lstlisting}

\clearpage

% =====================================================================
% 6. Reproducción
% =====================================================================
\section{Reproducción de la entrega}

\subsection{Requisitos}

\begin{itemize}
\item Python 3.12.10 (compatible con 3.10 o superior).
\item SQLite 3.49.1 incluido en Python; sin servidor ni credenciales.
\item Dependencias de \code{requirements.txt}: FastAPI, Uvicorn, pytest, httpx
  y coverage.
\end{itemize}

\subsection{Instalación y ejecución}

\begin{lstlisting}
python -m venv .venv
.\.venv\Scripts\Activate.ps1
python -m pip install -r requirements.txt

# Base de la API con el estado inicial conocido
python -m persistencia.inicializar_db

# Servidor local
uvicorn api.main:app --reload --port 8000
\end{lstlisting}

\begin{itemize}
\item Base URL: \url{http://127.0.0.1:8000/api/v1}
\item Documentación interactiva: \url{http://127.0.0.1:8000/docs}
\item OpenAPI: \url{http://127.0.0.1:8000/openapi.json}
\end{itemize}

{\footnotesize
\begin{tabularx}{\linewidth}{@{}V{0.6}V{1.1}V{1.3}@{}}
\toprule
\thead{Variable} & \thead{Valor por defecto} & \thead{Uso} \\ \midrule
\code{DB\_PATH} & \code{evidencias/api.db} & Ruta de la base SQLite. \\
\code{CORS\_ORIGINS} & \code{localhost:5173}, \code{127.0.0.1:5173} y
\code{localhost:3000} & Orígenes permitidos para la interfaz. \\
\bottomrule
\end{tabularx}
}

\subsection{Estructura relevante}

\begin{lstlisting}
proyecto_arq_sis_mascotas/
|-- api/                  # rutas, contratos y dependencias
|-- servicios/            # RF-05 y RF-34
|-- persistencia/         # esquema, seed, repositorio e inicializador
|-- pruebas/              # 39 pruebas y scripts de evidencia
|-- docs/
|   |-- contratos_api.md          # contratos completos
|   `-- informe_semanas_13_14.md  # informe en Markdown
|-- evidencias/           # pruebas, cobertura y desempeno
|-- plantilla/            # documentos formales
`-- README.md
\end{lstlisting}

\subsection{Evidencias}

{\footnotesize
\begin{tabularx}{\linewidth}{@{}V{0.7}V{1.3}@{}}
\toprule
\thead{Evidencia} & \thead{Contenido} \\ \midrule
\code{evidencias/salida\_pruebas.txt} & 12/12 pruebas de persistencia
(semana 13). \\
\code{evidencias/salida\_pruebas\_semana14.txt} & 39/39 pruebas de la
semana 14. \\
\code{evidencias/cobertura.txt} & Cobertura de servicios (RNF-08). \\
\code{evidencias/rnf\_desempeno.txt} & 20 consultas sobre 1.000 fichas
(RNF-05). \\
\code{evidencias/demo\_api.txt} & Peticiones y respuestas HTTP reales de ambos
RF. \\
\code{evidencias/prueba.db} & Base reproducible de la semana 13. \\
\bottomrule
\end{tabularx}
}

\clearpage

% =====================================================================
% 7. Checklist
% =====================================================================
\section{Checklist de entrega}

{\footnotesize
\begin{tabularx}{\linewidth}{@{}V{1.9}Z{0.35}V{0.75}@{}}
\toprule
\thead{Criterio} & \thead{Cumple} & \thead{Evidencia} \\ \midrule
Los dos RF mantienen la prioridad y el alcance de la semana 13 & \textbf{Sí} &
Sección 2 y \code{plantilla\_13\_1.pdf}. \\
Ambos RF tienen servicios y endpoints identificables & \textbf{Sí} &
\code{servicios/} y \code{api/main.py}. \\
Los contratos describen éxito, entrada inválida y errores & \textbf{Sí} &
\code{docs/contratos\_api.md} y \code{/docs}. \\
Hay resultado válido y rechazo/ausencia por cada RF & \textbf{Sí} &
Sección 4.6. \\
Se verifica la persistencia real y la ausencia de escrituras indebidas &
\textbf{Sí} & Pruebas con SQLite y conteos antes/después. \\
El SQL queda en persistencia y las reglas en servicio/dominio & \textbf{Sí} &
Sección 4.1. \\
Cada RF se vincula con sus RNF y atributos ISO & \textbf{Sí} & Sección 5. \\
Los RNF previstos tienen medidas y evidencias & \textbf{Sí} &
RNF-04, 05, 08 y 10; RNF-07 pendiente. \\
Se distinguen los incumplimientos y pendientes de la semana 3 & \textbf{Sí} &
Sección 8. \\
El README permite ejecutar y repetir las pruebas & \textbf{Sí} &
Sección 6 y \code{README.md}. \\
\bottomrule
\end{tabularx}
}

% =====================================================================
% 8. Pendientes
% =====================================================================
\section{Pendientes y relevo a la semana 3}

\textbf{Implementado y verificado}

\begin{itemize}
\item Persistencia completa de RF-05 y RF-34 con integridad referencial.
\item Servicios de aplicación y API REST con cuatro endpoints.
\item Contratos documentados y OpenAPI generado por FastAPI.
\item 39 pruebas funcionales, de integridad y de servicios aislados.
\item RNF-04, RNF-05 (secuencial), RNF-08 y RNF-10 con evidencia reproducible.
\item Demostración HTTP real de ambos RF.
\end{itemize}

\textbf{Pendiente para la semana 3}

\begin{itemize}
\item Autenticación, roles y reemplazo de \code{refugio\_id} explícito por el
  usuario autenticado.
\item Pantalla de ficha y catálogo (React/Vite) contra esta API.
\item Prueba de concurrencia de RNF-05 (50 usuarios) y RNF-07 (2.000 fichas,
  5.000 adoptantes).
\item Migración y verificación en PostgreSQL (ADR-03/ADR-10).
\item Almacenamiento binario de fotos (hoy se guarda la ruta).
\item Rechazo automático de solicitudes pendientes al dar de baja (RN-15),
  fuera de este corte.
\end{itemize}

% =====================================================================
% 9. Trazabilidad
% =====================================================================
\section{Trazabilidad}

{\footnotesize
\begin{tabularx}{\linewidth}{@{}V{0.6}V{1.4}@{}}
\toprule
\thead{Elemento} & \thead{Dónde se evidencia} \\ \midrule
RF-05 & \code{servicios/servicio\_mascotas.py}, \code{api/main.py}
(POST/PUT/DELETE) y \code{pruebas/test\_rf05\_api.py}. \\
RF-34 & \code{servicios/servicio\_catalogo.py}, \code{api/main.py} (GET) y
\code{pruebas/test\_rf34\_api.py}. \\
Contratos & \code{docs/contratos\_api.md}, \code{/docs} y
\code{/openapi.json}. \\
RNF-04 & \code{test\_catalogo\_sin\_datos\_personales\_rnf04}. \\
RNF-05 & \code{pruebas/rnf\_desempeno.py} y
\code{evidencias/rnf\_desempeno.txt}. \\
RNF-08 & \code{pruebas/test\_servicios\_aislados.py} y
\code{evidencias/cobertura.txt}. \\
RNF-10 & \code{pruebas/test\_rnf10\_integridad.py} y
\code{evidencias/salida\_pruebas\_semana14.txt}. \\
RNF-07 & Pendiente de la semana 3. \\
Persistencia & \code{persistencia/schema.sql},
\code{persistencia/repositorio.py} y \code{evidencias/salida\_pruebas.txt}. \\
Informe en Markdown & \code{docs/informe\_semanas\_13\_14.md}. \\
Documento semana 13 & \code{plantilla/plantilla\_13\_1.tex} y
\code{plantilla/plantilla\_13\_1.pdf}. \\
\bottomrule
\end{tabularx}
}

\section*{Cierre}

La semana 13 dejó una base reproducible y verificada; la semana 14 convirtió esa
persistencia en servicios y una API con contratos estables. Las dos semanas
mantienen el mismo alcance: publicar fichas (RF-05) y consultarlas (RF-34). La
calidad se concentró en integridad, comportamiento temporal, capacidad de prueba
y confidencialidad. Lo que falta ---autenticación, interfaz, concurrencia y
PostgreSQL--- está identificado y se retoma en la semana 3.

\end{document}
